\documentclass{article}

% if you need to pass options to natbib, use, e.g.:
%     \PassOptionsToPackage{numbers, compress}{natbib}
% before loading neurips_2021

% ready for submission
\usepackage[preprint]{neurips_2021}

% to compile a preprint version, e.g., for submission to arXiv, add add the
% [preprint] option:
%     \usepackage[preprint]{neurips_2021}

% to compile a camera-ready version, add the [final] option, e.g.:
%     \usepackage[final]{neurips_2021}

% to avoid loading the natbib package, add option nonatbib:
%    \usepackage[nonatbib]{neurips_2021}

\usepackage[utf8]{inputenc} % allow utf-8 input
\usepackage[T1]{fontenc}    % use 8-bit T1 fonts
\usepackage[colorlinks=true]{hyperref}       % hyperlinks
\usepackage{url}            % simple URL typesetting
\usepackage{booktabs}       % professional-quality tables
\usepackage{amsfonts}       % blackboard math symbols
\usepackage{nicefrac}       % compact symbols for 1/2, etc.
\usepackage{microtype}      % microtypography
\usepackage{xcolor}         % colors
\usepackage{graphicx}       % Extended version of graphics
\usepackage{subcaption}

\newcommand{\todo}[1]{\textcolor{red}{\textbf{#1}}}

\graphicspath{ {./images/} }

\title{Is lottery fair?}

% The \author macro works with any number of authors. There are two commands
% used to separate the names and addresses of multiple authors: \And and \AND.
%
% Using \And between authors leaves it to LaTeX to determine where to break the
% lines. Using \AND forces a line break at that point. So, if LaTeX puts 3 of 4
% authors names on the first line, and the last on the second line, try using
% \AND instead of \And before the third author name.

\author{%
  Michal Horáček\\
  Matrikelnummer 6373382\\
  \texttt{michal.horacek@student.uni-tuebingen.de} \\
  \And
  Carson Zhang\\
  Matrikelnummer 6384481\\
  \texttt{carson.zhang@student.uni-tuebingen.de} \\
}

\begin{document}

\maketitle

\begin{abstract}
    We investigate the uniformity of various lotteries. 
    First, we test the null hypothesis that, for the German Lotto, the distribution of the minimum distance between
    lottery numbers, $d$, is as it would be in a uniformly sampled lottery. The $\chi^2_6$ goodness-of-fit
    test resulted in a p-value of $0.4989$.
    Furthermore, to evaluate the randomness of the random number generators underlying different lotteries,
    we apply the Diehard battery of tests \todo{with the following results}.
\end{abstract}

\section{Introduction}
\label{sec:introduction}
Lottery/Hazard business is big many. Games such as lottery and dice are (probably?) one of the first application of randomness in human culture.
But for the same time organizers have cheated in lotteries, skewing the uniform distribution of drawn numbers expected by common sense for their
own monetary gain. We investigate the distribution of answers for the German Lotto lottery drawn from 1955 to the present day.

In the second part of this paper, we explore the "quality" of random numbers generated by lottery draws via the Diehard battery of tests~\cite{diehard}. These
tests were developed in the 1990's by George Marsaglia and are still in use today to test random number generators~\cite{modernRNG}.


\section{Dataset}
\label{sec:dataset}
Section \ref{sec:methods} investigates numbers drawn from the German lottery Lotto, from 1955 onwards. 
We have chosen this dataset because it is relevant to us and readers of this analysis as current residents of Germany, 
and it contains almost 70 years of data, which we believ is enough data to perform a meaningful analysis.

Furthermore, Drakakis, Taylor, and Rickard did not apply their tests to this data, so 
performing the test on the German Lotto data is a meaningful new result.

The dataset has been compiled by Johannes Friedrich, 
a software developer who has made the data publicly available via 
\href{https://github.com/JohannesFriedrich/LottoNumberArchive}{his Github repository}.

\todo{Perform verification of the correctness of the dataset by sampling from it and manually inspecting the data.}

\todo{Is there a dataset for a rigged lottery?} \href{https://notebook.community/JesseScott/Lotto649/lotto}{This} looks sketchy as hell.

However even 70 years of Lotto numbers is not sufficient to produce enough data for the diehard tests. Their author recommends
using 10 to 12 MB of random bits, which is substantially more than 28.4 KiB of Lotto numbers. Therefore we add winning numbers
from other lotteries to our dataset. By data volume, our data is dominated by the New York Quick Draw lottery, because it has
been drawn every 4 minutes for the last decade, adding up to to 25,663,100 numbers at our collection of the dataset.
Washington DC's Keno is quite similar, reaching 2,299,563 numbers over 3 years of existence. Our dataset is completed by
the "other" category, consisting of 1,917,890 winning numbers combined from lotteries drawn less often worldwide.

\begin{figure}
    \centering
    \includegraphics[width=\textwidth]{dataset_composition.pdf}
    \caption{Origin of numbers making up our dataset.}
    \label{fig:dataset}
\end{figure}

The "other" dataset comes from European or english-speaking countries. The reason is twofold - we have been looking
for them with English search queries and these nations are more likely to subscribe to ideas like open data and thus
offer such dataset as csv files. The European nationality stands for European transnational lotteries Eurojackpot and Euromillions.

\begin{figure}
    \centering
    \includegraphics[width=\textwidth]{other_composition.pdf}
    \caption{National composition of the "other" numbers.}
\end{figure}

In total, we have gathered winning numbers from 37 different lotteries.


\section{Methods}
\label{sec:methods}
We conducted a hypothesis test. Null hypothesis $H_0$ assumes that lottery numbers are sampled uniformly,
alternative hypothesis $H_A$ assumes the opposite. We used Pearson's $\chi^2$ test, a commonly recommended
test for the probabilities of observing categorical data (\todo{sOuRcE?}). Furthermore, it is the same type
of test performed by Drakakis et al.~\cite{Drakakis}, which allows us to check our (and their) results.
% We can also satisfy the assumptions of the $\chi^2$ test, and 
% follow the common rules-of-thumb for using the test. \todo{Motivate the choice of this particular test.}

\subsection{Minimum-distance statistic}

To create the frequency tables for the $\chi^2 test$, we used the minimum distance statistic $d$ used by Drakakis et al.
This statistic is useful for detecting human tampering because it is known that humans usually do a
poor job of imitating the random choices that occur under the null hypothesis. Boland and Pawitan \todo{cite} showed
that when humans are asked to sample $m$ integers from the set ${1,...,n}$, the value
of this statistic is greater on average than it is under uniform sampling. In particular,
they observed that the true probability of $d = 1$ is greater than $0.5$, a result that 
is highly unintuitive to humans: humans tend to underestimate how likely it is that
two consecutive numbers are picked. In particular, Boland and Pawitan
computed a $\chi^2$ goodness-of-fit statistic of the human-produced $d$ against the uniformly produced $d$
, which if they had conducted a hypothesis test, would have yielded a 
p-value of $< 0.0000001$. 

Therefore, a low p-value could be consistent with human tampering, especially if we
observe an unusually small number of lottery drawings where $d = 1$, as Drakakis et al.
did in the French lottery. 

\subsection{The test statistic}

The $\chi^2$ goodness-of-fit test tests how well an expected discrete distribution fits
an observed discrete distribution. We computed the expected distribution using the following formula proved by Drakakis (2007) \todo{cite}.

Let the lottery game be a sample of $m$ integers drawn from the integers $1,...,n$. 
(In the German Lotto, $n = 49$ and $m = 6$.) Let  $r_1,...,r_m$ be the numbers drawn in the sample of size $m$.
The minimum distance
\begin{equation}
    d = min_{1 \leq i < j \leq m} | r_j - r_i |
\end{equation}
has the following distribution
\begin{equation}
     P(d < k) = 1 - \frac{{n - (k - 1)(m - 1) \choose m}}{{n \choose m}}, \quad k = 1,..., \left\lfloor \frac{n - 1}{m - 1} \right\rfloor.
\end{equation}

We define $\alpha = 0.05$ as our significance level because it is popular and Krakakis used
the same significance level.

To prepare the German Lotto dataset for the hypothesis test, we computed $d$, the minimum distance between winning numbers, for each 
lottery day. We counted the frequencies of each value of $d$, and combined the counts of the 
two largest possible values, $7$ and $8$. This is the same type of data preparation that was performed in the paper,
and we did it to satisfy the common rule of thumb for usage of the 
$\chi^2$ goodness-of-fit test: the expected frequency in each bin must be $\geq 5$. This prepration yielded the dataset in Table \ref{tab:frequencies}.

\begin{table}
    \centering
    \caption{Frequencies of $d$ statistic for German Lotto}
    \begin{tabular}[t]{lrr}
        \toprule
        d & Expected frequency & Actual frequency\\
        \midrule
        1 & 2310.595965 & 2383\\
        2 & 1266.759926 & 1247\\
        3 & 639.888057 & 615\\
        4 & 290.255446 & 273\\
        5 & 113.589766 & 105\\
        \addlinespace
        6 & 35.857667 & 35\\
        7 and 8 & 9.025145 & 8\\
        \bottomrule
    \end{tabular}
    \label{tab:frequencies}
\end{table}

We estimated no parameters from the data, so we use
$p - 1 = 7 - 1 = 6$ degrees of freedom for the $\chi^2$ statistic.


% \section{Results}
% \label{sec:results}
% We have applied the tests and found that we are indeed very smart.

\section{Diehard tests}
\label{sec:diehard}
Fairness entails more than the question whether are lottery numbers from the expected distribution. For instance, the Kolmogorov-Smirnov
test used in the first part of this paper does not concern itself with the order the numbers are drawn. However if we saw a lottery whose
numbers were always drawn in a descending sequence, for example, we would become suspicious.

Thus a more comprehensive test is clearly required to establish a more detailed answer to our question. We approach this problem by reformulating
lottery as a process producing a stream of (supposedly) random numbers, which themselves are simply bit sequences. Under this formulation,
we can deploy standard statistical tests developed for testing random number generators: we have a file of one and zero bits and wish to investigate
if its bits are correlated, repeating with a period or other quantities undesirable for randomness.

A number of these test suites has been developed over time. Donald Knuth presented an initial set of empirical tests in the second volume of his computer science bible
The Art of Computer Programming in 1969. Many general cryptography textbooks such as \href{http://www.cacr.math.uwaterloo.ca/hac/}{Handbook of Applied Cryptography} 
or \href{http://www.wisdom.weizmann.ac.il/~oded/foc-vol1.html}{Foundations of Cryptography} contain multiple tests of their own. The american National Institute 
of Standards \& Technology has published a \href{https://nvlpubs.nist.gov/nistpubs/legacy/sp/nistspecialpublication800-22r1a.pdf}{guideline} discussing this matter too.

We decided to use the Diehard battery of tests, which was developed by the american statistician George Marsaglia in the nineties. While this package used
to be quite popular in its day, it has been superceded today by other suites, including its derivatives such as Dieharder or TestU01. \todo{Why did we pick this one?}

\begin{enumerate}
    \item About Diehard.
    \item The following part is divided into several section. In section 1, we discuss data gathering, processing and general creation of input files for Diehard. Special
    attention is given to the problem of attaining input file of sufficient size and the asymetric requirements of individual Diehard tests.
    \item Part 2 highlights result on Diehard suite, interprets them and compares them against commonly used PRNGs.
    \item Section 3 clarifies the shortcomings of above approach.
\end{enumerate}


\section{Conclusion}
\label{sec:conclusion}
We analyzed the uniformity of lottery results in two major ways. We performed a hypothesis test of the distribution of the minimum distance $d$, a statistic that is known to be poorly imitated by humans. Based on the results for the German Lotto and NY Quick Draw, there is no strong evidence for human tampering, but it is unknown how well this analysis may detect other forms of unfairness, and it is unknown how well this would inform a player's choices.

% The Diehard analysis was greatly limited by insufficient amount of data and in some cases unsuitable format in which the data was reported in (ascending-order draws). Out of the three datasets, we can only make a semi-conclusive statement about DC Keno. Here the evidence points towards the validity of the null hypothesis; that is, the randomness and therefore fairness of the lottery.

The Diehard analysis was greatly limited by having too little data, the without-replacement nature of lottery sampling, and some obviously non-random (sorted) data. Despite these concerns, we believe that the performance of DC Keno on the Diehard tests suggests the fairness of this lottery.


\newpage

\bibliography{sources.bib}
\bibliographystyle{unsrt}

\end{document}
